% !TeX root = ../main.tex


%%                  %%
%%                  %%
%%  16 de outubro   %%
%%                  %%
%%                  %%
\entrada{16 de outubro de 2026}{Pensar a comunicação}

Esta pesquisa começa no dia 2 de abril de 2024 -- data que marca o dia internacional de conscientização do autismo -- quando vi numa rede social uma publicação condenando os memes com características de transtornos psíquicos compartilhados por quem vive com estes.

Sendo eu uma pessoa autista, me incomodei consideravelmente com aquela provocação naquele dia específico. Memificar nossas características em comum é um jeito de falarmos de nossos corpos por vias que operam desvios frente ao utilitarismo do vocabulário diagnóstico; é um jeito de constituir com nossos corpos identidades mediadas positivamente, criativamente, e não pelo intermédio do defeito. Na lembrança daquela ocasião, enxergo com os olhos de hoje os erros que Latour atribui aos pós-modernos: tenta-se desconstruir os diagnosticados apagando junto às doenças as diferenças mediadas hoje pelos diagnósticos; a crítica despatologizante de nada ajuda aos patologizados podando suas diferenças\notaarquivo{notas/latour-posmod.tex}.

Decidi que escreveria um texto respondendo à postagem, buscando reintegrar as diferenças ao discurso sobre o corpo, e comecei a pesquisa que articularia tal escrita. Logo a proposta evoluiu para um trabalho maior, que almejei transformar num projeto de iniciação científica (o qual não chegou a ser elaborado): eu queria falar de autismo nas veredas entre as ciências médicas --- patologizantes --- e as vertentes das humanidades críticas à biologização da vida. Eu queria manter as diferenças sem enquadrá-las como defeitos, assim como trazer os problemas do paradigma diagnóstico sem apagar as diferenças; busquei, então, resolver a controvérsia adicionando-lhe uma dimensão: os quase-sujeitos em questão.

Latour descreve a Constituição dos modernos como um princípio de proliferação de quase-objetos -- mediadores que são deixados pelo caminho da estabilização de sujeitos e objetos. Digo que somos enquadrados como quase-sujeitos na controvérsia em questão porque, de um ao outro lado, traduzimos circunstancialmente uma polarização ética e teórica na qual nossas diferenças específicas são tolhidas de qualquer valor, quando não da própria voz. Aqueles que patologizam as diferenças, ao menos têm o cuidado de reconhecê-las; enquanto os que se negam aos rótulos diagnósticos, vão tecendo diagnósticos para o mundo, apagando as diferenças de habituação entre corpos distintos\notaarquivo{notas/mundo-doente.tex}.

A questão foi ficando ainda mais interessante pois incorporaria também uma dimensão política. A diferenciação pelo enquadramento da doença ou da deficiência medeia o acesso a certos direitos, além de articular discursos que mobilizam solidariedades. Como fazer, portanto, para trocar o figurino de defeito de nossas diferenças por outro -- de valência positiva -- sem acarretar desatenção aos modos específicos como elas nos proporcionam existir?

Busquei a orientação do professor Pedro Peixoto por ter sido o único em meu instituto em cujo currículo encontrei alguma menção ao autismo. Seu trabalho no campo dos estudos sociais das ciências e técnicas o coloca em companhia frequente com o pensamento de Bruno Latour -- intelectual que eu até então desconhecia. Pensando com Latour e Peixoto, compreendi que faltava uma dimensão ao enquadramento moderno da comunicação: para que as diferenças patologizadas sejam positivadas, o significado de uma mensagem deveria ser o resultado, sempre incerto, de uma performance, e não um atributo rígido da própria mensagem.

O significado é um acontecimento no emissor, e é outro no receptor --- que constrói um significado à parte, articulável em alguma medida com o primeiro. Tal é a medida da confluência comunicacional\notaarquivo{notas/confluencia_comunicacional.tex}. Com isso minha pesquisa tomou outro rumo: o de compor enquadramentos para a comunicação e a ciência (enquanto processo genérico pelo qual se viria a ter ciência de algo) de modo a agregar o ruído como parte da mensagem, e não como um sinal que deve ser descartado a fim de melhor significá-la.

%%                  %%
%%                  %%
%%  17 de outubro   %%
%%                  %%
%%                  %%
\entrada{17 de outubro de 2026}{Pensar o antropoceno}

O dossiê para o qual escrevo pede pensar o antropoceno, e minha primeira entrada neste diário de escrita -- escrita que é o próprio diário\notaarquivo{notas/escrita_experimental} -- não parece ter se aproximado nem introduzido o tema. Penso, então, em dois caminhos que realizam este contato: um que  segue a cronologia linear que descrevi anteriormente; outro que rompe com ela e espirala a partir da inadequação dos enquadramentos modernos à comunicação e à ciência.

Seguindo agora pelo segundo, proponho pensar o Antropoceno a partir dessa inadequação. Encontro mais agência diante das emergências naturais e sociais do presente nomeando um \enquote{Antropoceno} --- com aspas --- que revela em seu nome a elipse dos quase-objetos quase-sujeitos, e na presença das aspas um estrangeirismo do termo. Demarcar uma época geológica pelos impactos humanos sobre o sistema planetário\notaarquivo{notas/da-maquina-a-vapor-ao-plutonio.tex} figura seres não-humanos -- como montanhas, rios, oceanos, florestas, a atmosfera, o clima, entre tantos outros -- como objetos; como se, de fato, tivessem sido meros recursos, e os ditos impactos humanos resultassem da agência humana, e não de uma associação de agências humanas e não-humanas; já a demarcação de um nome com as aspas ao mesmo tempo que nos livra de falar como se fôssemos antropocênicos\notaarquivo{notas/antropocenicos.tex}, reintegra a dimensão que estes retiram da mensagem; a visão simultânea de ambas as demarcações proporciona agir-comunicar-com quase-objetos quase-sujeitos  mais articuladamente.

Foquemos agora na primeira demarcação: a nomenclatura do \enquote{Antropoceno}. Em sua escolha pelos antropocênicos revela-se a autovisão destes, assim como seu caráter precipitado. Enquadram-se como uma força geológica, mas o fazem a partir da observação de efeitos geológicos não previstos nem desejados da aplicação desta força; enquadram-se como uma força,  mas desgovernada, imprevisível; o fazem pois não conseguem enxergar agenciamento e delegação entre os seres não-humanos, porque se precipitam e tomam como objetos seres que são afetados e se movem como sujeitos, pois são incapazes de enquadrar não-humanos enquanto tais.

Se a atmosfera for figurada como sujeito, demarcar  transformações climáticas nocivas aos humanos como  antropogênicas seria como atribuir aos senhores de escravos mortos na Revolução Haitiana a autoria maior deste evento, e não àqueles humanos que reagem à escravidão alçados por seus próprios interesses. Ao retirar-lhes a dignidade, cegos aos interesses não-humanos, os humanos convidam-nos às animosidades. Nos quilombos, conforme ensina Nego Bispo, compartilha-se a dignidade com o outro e com os não-humanos, convidando-se todos às confluências\notaarquivo{notas/bispo-confluencias.tex}.

Ainda que seja para denunciar os efeitos adversos das agências humanas para a habitabilidade planetária a estes agentes, a nomenclatura do \enquote{Antropoceno} se fundamenta no mesmo monoteísmo que os orienta.

%%                  %%
%%                  %%
%%  18 de outubro   %%
%%                  %%
%%                  %%
\entrada{18 de outubro de 2026}{Retornar ao que foi comunicado}

Ontem, parei a escrita logo após uma menção a Nego Bispo. Nem sequer terminei de escrever o último parágrafo\notaarquivo{notas/paragrafo.tex}. Comecei a escrevê-lo  mas  interrompi a escrita pois o mestre quilombola chamava a uma releitura de "A terra dá, a terra quer"\notaarquivo{notas/refs/a-terra-da.tex} antes de retomar a composição do texto com uma presença maior de Bispo.

Leitura inspiradora e potente, cuja leveza espanta. Concluí o reencontro em  menos de 24 horas. Decidi em meio ao acontecimento presentear o livro a duas pessoas às quais suas palavras me levaram todo o tempo: uma, pela confluência com a ancestralidade; outra, pela desejo de viver a vida no mato.

Me intrigou notar como Bispo constitui riquíssimas encruzilhadas. Dos modos de vida rurais com os saberes ancestrais; mas também entre os modos de se relacionar com a terra e os modos de se relacionar com outros povos humanos. Esta segunda encruza me levou a pensar em dois textos: um do antropólogo francês Marcel Mauss\notaarquivo{notas/refs/var-sazonais.tex}, outro da antropóloga abrasileirada Manuela Carneiro da Cunha\notaarquivo{notas/refs/cultura.tex}.

Com seu texto, Mauss me deixou fascinado com os povos Inuit (nomenclatura contracolonial daqueles que o francês chamou de esquimós\notaarquivo{notas/inuit.tex}). O estudo, elaborado conjuntamente com Henri Beuchat, traduz de seus modos de vida uma coesão ancestral que atravessa inúmeras comunidades, cada qual confluindo com a terra, com o tempo e entre si. Reler Antônio Bispo nos ensinando os modos de vida quilombolas fez ressoar em mim a mesma ancestralidade confluente entre os Inuit.

Vamos pensar, por exemplo, na sensibilidade geracional deste povo e dos quilombolas. Os primeiros mantêm  em suas coletividades  um mesmo conjunto de nomes pessoais, sendo as  crianças batizadas e encarnadas com o nome dos últimos ancestrais  falecidos. A geração neta revive a geração avó; para Bispo e seus conterrâneos, começo-meio-começo: geração avó, geração mãe, geração neta\notaarquivo{notas/mauss-bispo-geracoes.tex}.

Na época das secas, os quilombolas se refugiam nas sombras das casas de farinha e dos engenhos -- é o período do ano em que a vida coletiva é mais efervescente. Todos se encontram nas casas, trocam o que têm pelo que necessitam, ajudam com o trabalho caso não tenham nada, festejam, conversam; acima de tudo, trocam saberes, afetos e vivências, se envolvem entre si. No restante do ano, há sombras no mato e cada um vai trabalhar em sua própria roça e a vida comunitária fica mais restrita ao núcleo familiar e à vizinhança\notaarquivo{notas/bispo-casas-de-farinha.tex}.

Bispo relata uma vida social que respira junto com a terra nos quilombos; movimento que se encontra também na tradução por Mauss dos modos de vida Inuit; há em sua descrição duas configurações sociais que são igualmente determinadas pelas estações do ano: no verão, os diferentes núcleos familiares vivem isolados cada um em um abrigo próprio construído no início da estação, e pouco interagem com a vizinhança no cotidiano; contudo, no inverno, as famílias se unem em assentamentos e vivem todas num mesmo abrigo por assentamento, período em que as relações com a terra são mais amenas, e as do coletivo são intensas, havendo muitas festas, brincadeiras e compartilhamento de recursos durante o período. O francês relata entre os Inuit uma respiração ainda mais acentuada do que aquela narrada por Bispo; ambas, no entanto, em profunda confluência com a terra.

Confluência que se avista também nas construções de cada povo. Ambos conversam intimamente com seu território, tecendo relações contratuais de profundo respeito e envolvimento. As construções Inuit variam, assim como as quilombolas, de acordo com sua localização; utilizam o que o território local tem a oferecer para construir casas e ferramentas atentas às condições de existência impostas e disponíveis localmente\notaarquivo{notas/mauss-bispo-construcoes.tex}.

Penso com as palavras de Bispo e Mauss em um saber ancestral comum a Inuits e quilombolas: o saber de confluir com a terra e todos os seus compartilhantes. E a ancestralidade me leva então para as últimas palavras de Ailton Krenak em \enquote{Futuro Ancestral} -- livro que rima na poesia, no espírito e na carne com o de Bispo, e termina dizendo: \enquote{o que nossas crianças aprendem desde cedo é a colocar o coração no ritmo da terra.}\notaarquivo{notas/refs/futuro-ancestral.tex} Ouso pensar com Krenak que este saber ancestral indígena, que conflui com a ancestralidade de seus compartilhantes do extremo norte continental -- os Inuit -- e com a de seus compartilhantes quilombolas, constitui um saber ancestral de todos os povos da terra, de todos os seus compartilhantes, humanos ou não, e que se perdeu diante do monoteísmo cosmofóbico europeu-modernizante-ocidental.

Retomando -- enfim -- o tema deste estudo-carta-diário não-antropocênico sobre como pensar o \enquote{Antropoceno}, penso que a Constituição moderna tecida por Latour rima com minha proposição dos antropocênicos e sua patologia comunicacional; mas rima ainda com um terceiro diagnóstico, o da cosmofobia enquanto patologia da ancestralidade de um povo que desaprendeu como respirar e colocar o coração no ritmo da terra, que se esqueceu deste saber e o pratica muito precariamente na vida higienizada e humanizada, desconectada da terra e de seus compartilhantes, que se enraizou nas cidades\notaarquivo{notas/cosmofobia.tex}.

%%                  %%
%%                  %%
%%  19 de outubro   %%
%%                  %%
%%                  %%
\entrada{19 de outubro de 2026}{Semeaduras}

Nego Bispo levou meu pensamento a semeaduras; com elas em mente, comecei uma releitura do texto supracitado de Manuela Carneiro da Cunha. Depois da rima, tripla, entre Bispo e Latour, trago outra, agora entre o quilombola e a antropóloga. Quando o leio dizendo que a agricultura quilombola não é mercadoria, remeto ao relato da antropóloga sobre a fala de um velho chefe yawanawa de que \enquote{\textit{Honi} não é cultura!}\notaarquivo{notas/manuela-honi.tex}

As articulações entre a crônica trans-mitológica de Nego Bispo e a crônica trans-étnica de Manuela Carneiro já se iluminavam desde as primeiras páginas de \enquote{A terra dá, a terra quer}. \enquote{Precisamos transformar as armas dos inimigos em defesa}\notaarquivo{notas/refs/a-terra-da-defesa.tex}, dizia um mestre de Bispo, que o quilombola cita para introduzir a semeadura de palavras e sentidos na língua dos humanos como proteção contra o colonialismo destes. Aqui se tem uma tecnologia ancestral de defesa interétnica ao avesso daquela descrita por Carneiro.

A antropóloga descreve o uso do estrangeirismo \enquote{cultura} por povos indígenas amazônicos, analisando mais especificamente a ocasião\notaarquivo{notas/manuela-ocasiao.tex} de uma negociação de direitos intelectuais sobre o uso em \textit{kampô} --  uma toxina encontrada na pele de certas rãs. Embora possuam em suas línguas palavras para nomear o complexo de seus costumes e modos de vida, a palavra \enquote{cultura} foi incorporada aos vocabulários de tais línguas; quando o velho chefe diz que \textit{honi} não é \enquote{cultura}, ele não nega que o chá faça parte das práticas que organizam e orientam a vida de seu povo. O estrangeirismo nomeia a \enquote{cultura} que os humanos colonizadores vão negociar; com ele, o velho chefe bloqueia os efeitos de reflexividade semântica que a tradução para termos nativos traria.

Lembro de Latour dizendo que, diferentemente dos modernos, que interditavam o pensamento nos híbridos de natureza e cultura, os pré-modernos dedicavam-se a pensá-los, e impediam, assim, sua proliferação desenfreada\notaarquivo{notas/latour-mistura.tex}. Isso é bem visível aqui, pois \textit{honi} não é \enquote{cultura} -- mesmo sendo claramente um componente da cultura daquele chefe -- porque não atende às relações sociais articuladas mediante o estrangeirismo. Objetos e sujeitos, natureza e sociedade, estão misturados nas palavras, de um modo que colabora para a manutenção de sua cultura sem enrijecê-la.

Na ocasião narrada pela antropóloga, o \textit{kampô} é \enquote{cultura}, mas \textit{honi} não é. Isto lhe causa intriga. A diferenciação é imunização contra o colonialismo: porque nos imaginários humanos, os direitos intelectuais sobre algo só podem ser individuais, coletivos ou públicos; nas relações com os povos da terra, isso restringe as negociações às formatações presentes nos imaginários metropolitanos; o uso de estrangeirismos em tais contextos interétnicos, além de melhor demarcar semanticamente os termos utilizados, protege as relações sociais encarnadas na língua nativa das contingências dos imaginários estrangeiros.

\textit{Honi} não é \enquote{cultura} porque as relações sociais ancoradas neste termo estão limitadas pelo imaginário ocidental, e se articulam ali a uma noção de propriedade coletiva. Organizações indígenas aprenderam a falar a língua do colonizador para se protegerem de sua exploração, e a \enquote{cultura} é uma de suas principais defesas. \textit{Honi} não é \enquote{cultura} porque, apesar de fazer parte da cultura nativa, não se articula a uma noção de propriedade coletiva, mas a um regime de conhecimento estranho ao imaginário ocidental, protegido de sua dominação mediante a diferenciação vocabular.

Se Nego Bispo semeia palavras e sentidos no vocabulário humano para proteger os modos de vida quilombola, Carneiro da Cunha relata um uso invertido do mesmo mecanismo. Semeando o vocabulário nativo com estrangeirismos metropolitanos, o velho chefe imunizava \textit{honi} contra sementes colonizadoras. Parafraseando Donna Haraway, importam muito mesmo quais nomes nomeiam as coisas\notaarquivo{notas/refs/haraway-e-os-nomes.tex}.

Como eu havia dito ontem, encontrei em \enquote{A terra dá, a terra quer} uma encruzilhada entre os modos como os povos da terra se relacionam com ela e como se relacionam com outros humanos. Bispo articulou em mim o primeiro elemento com o relato de Mauss sobre os Inuit, e o ,segundo com o de Manuela sobre a \enquote{cultura} entre povos indígenas amazônicos. Em ambos, encontro uma sensibilidade fina às circunstâncias locais -- ao que se repete e ao que se diferencia em cada situação -- que me deixa tentado a situar como o saber ancestral mais fundamental a seguinte proposição: sempre escutar e aprender a escutar as diferenças --- as ancestrais e as emergentes\notaarquivo{notas/dif-ancestral-emergente.tex}.

%%                  %%
%%                  %%
%%  20 de outubro   %%
%%                  %%
%%                  %%
\entrada{20 de outubro de 2026}{Uma patologia da comunicação}

Latour escreveu um ensaio sobre falas do corpo\notaarquivo{notas/refs/how-to-talk.tex} com o qual comecei a pensar a proposta de Cecinepá como uma reação ao antropoceno. Trata-se de um texto muito importante para minha pesquisa desde a primeira vez que o li. Cecinepá é minha contribuição para as proposições ali apresentadas de composição progressiva e articulada de um mundo comum --- um multiverso --- e de ciências mais inter-essantes e inter-essadas.

Hoje, quero trabalhar a noção de que o antropoceno é uma policrise derivada de uma patologia generalizada da comunicação: a patologia de não enxergar que, entre o sujeito e o objeto, há uma forma que articula a emergência polarizada de ambos. O ensaio de Latour trabalha muito esta relação -- em minhas palavras -- entre sujeito, forma e objeto. Cecinepá é o componente mitológico da imunização que componho a fim de nos protegermos de tal patologia.

Latour assume o corpo como uma interface em constante processo de aprendizagem -- portanto, de transformação. Através desse processo, o corpo se torna progressivamente sensível a mais diferenças encontradas no mundo. A aprendizagem cultiva verdadeiros órgãos de sentido no corpo aprendiz; num só ato, ela faz emergirem um novo meio sensorial e um novo mundo sensível; transformam-se sujeito e objeto de uma só vez, articulados entre si.

Para ilustrar o processo, o francês apresenta o caso de perfumistas aprendizes e seu treinamento com kits de odores para se tornarem narizes\notaarquivo{notas/nariz.tex}. O kit-de-odores é um artefato construído por um especialista ao ordenar um certo número de pares de fragrâncias desde as mais marcadamente diferenciadas até aqueles que apresentam os contrastes mais sutis. Após uma semana expondo-se continuamente à sequência do kit, os corpos aprendizes progressivamente constituem narizes sensíveis a diferenças olfativas que, até então, lhes passavam despercebidas, mudas; transforma-se, assim, o corpo-sujeito, com um novo órgão olfativo; transforma-se, também, o mundo-objeto que ele habita, recheado agora de novos odores.

O kit-de-odores e o processo de treinamento constituem a forma que proporciona a emergência deste corpo-mundo --- sujeito-objeto. O treinamento e o kit ordenam e propõem uma forma aos contrastes olfativos apresentados; articulam suas diferenças à forma da linearidade, da sequência linear, já sensível a corpos europeus-modernizantes-ocidentais. Não só aquelas diferenças mais sutis vão ficando mais marcantes por se situarem num lugar específico da sequência, mas os afetos que cada odor porventura invocar em cada corpo vão se alinhando também; esta sequência de odores e de contrastes é uma forma composta artificialmente que \guillemetleft ganha corpo\guillemetright ao ressoar uma forma linear já articulada pelo corpo aprendiz junto a uma sequência de diferenças sensíveis imperceptíveis a ele ecoando e formatando referenciais afetivos subjetivos, circunstanciais a cada corpo; o kit-de-odores, enquanto uma metaforma, proporciona a sensibilização de um corpo-mundo a novos elementos que passam a compor um corpo`-mundo` -- a dupla articulada sujeito-objeto se modifica quando exposta continuamente a uma forma-sujeito (os pares de odores) articulada a uma forma-objeto (a linearidade), até que emerjam aderências entre a forma-sujeito e uma segunda forma-objeto (os afetos subjetivos), interna ao corpo; com isso, a forma-sujeito passa a ser localizada por duas formas-objeto, uma externamente proposta e outra internamente constituída; esta triangulação faz com que a forma-sujeito passe a ser ouvida como forma-objeto, de um modo que ainda verifica a consistência entre a articulação proposta pelo kit e o mundo original habitado por aquele corpo; a forma -- o treinamento com o kit-de-odores -- cultiva a emergência de um novo par corpo-mundo, ou sujeito-objeto -- ou, ainda melhor, sujeito-oriente.

Abolir o objeto em prol do oriente me soa bem por pelo menos dois grandes motivos. Primeiro, objeto remete a manipulação, enquanto oriente remete a orientação, localização, o máximo de confiança que podemos esperar dos seres, das diferenças no mundo; segundo, o oriente contrapõe o ocidente, coletivo articulado à patologia à qual busco a imunização. Encontro, então, um terceiro motivo: há honra em ser um oriente para o outro, não há em ser um objeto.

Assim, rearticulo: o kit-de-odores articula uma forma-sujeito a uma forma-oriente; através da exposição repetida a esta articulação, que chamarei de artifício, um corpo aprendiz progressivamente articula a forma-sujeito a uma segunda forma-oriente; aquela então se estabiliza enquanto forma-oriente\notaarquivo{notas/forma-oriente.tex}; o corpo-mundo se atualiza em corpo`-mundo`. O artifício \guillemetleft ganhou corpo\guillemetright. Através da exposição continuada a um corpo, o artifício -- uma metaforma composta por uma forma-sujeito articulada a uma forma-oriente -- se faz emergir como forma-oriente num corpo-mundo transformado; da metaforma germina, deste modo, um corpo-mundo transformado, a dualidade sujeito-oriente atualizada.

O sujeito é desorientador\notaarquivo{notas/sujeito.tex}, forma sem forma, inconsistente; quando não, imperceptível. A forma-sujeito só aparece carregada por uma forma-oriente que aponte: aqui está uma forma que você não pode prever, ou mesmo perceber, a não ser que eu a indique. A familiaridade adquirida com um sujeito pode articulá-lo a algum outro oriente e transformá-lo, enfim, numa forma-oriente.

A patologia antropocênica decorre de uma elisão da forma enquanto mediação criativa do sujeito e do mundo que o orienta. Formas são os híbridos, laboratórios, experimentos, escolas, universidades, os artifícios diversos que se proliferam e somem; são localizações perdidas; como se o mundo informasse de si mesmo. Perdemos a capacidade de comunicar -- de construir o comum -- porque perdemos as aspas, a capacidade de navegar e se articular entre diferentes regimes de conhecimento, narrações do mundo, mitologias, teorias, disciplinas, ciências, religiões, e por aí vai.

Nego Bispo denuncia o monoteísmo dos povos da cidade, e seus críticos acham que ele fala de religião\notaarquivo{notas/criticos-bispo.tex}. Monoteísmo é a cosmofobia, o medo do Outro, a incapacidade de se articular a diferentes referenciais cosmológicos; ou seja, a diferentes deuses. O monoteísmo é doença que impede a comunicação, e seu mecanismo pode ser descrito pela "Constituição dos modernos" transcrita por Latour.

Natureza transcendental, entretanto, agenciada como objeto para a transformação da realidade; sociedade imanente que se assume anterior ao próprio sujeito; um abismo separando natureza e sociedade; e um deus barrado, juiz nada-poderoso ao qual se recorre nos períodos de crises e conflitos morais. Esta é a forma da elipse patológica, que se encontra mais nitidamente no signo enquanto categoria da linguagem, da comunicação e da semiótica: um significante e um significado, um sujeito e um objeto, um ativo e um passivo, separados, hierarquizados, o sujeito-sociedade-significante, o objeto-natureza-significado, com um deus barrado se escondendo no meio. O signo vem servindo como um cabresto que dificulta a articulação dos antropocênicos entre diferentes referenciais cosmológicos -- ou regimes de saber, ou \enquote{culturas}, ou \enquote{verdades}, ou teorias, ou ciências, ... -- e daí se dificulta cumulativamente a comunicação enquanto composição de uma coletividade de diferenças em torno de um referencial comum, mas não único.

%%                  %%
%%                  %%
%%  21 de outubro   %%
%%                  %%
%%                  %%
\entrada{21 de outubro de 2026}{Forma}

Ando indeciso quanto à forma deste texto. Acho que eu gostaria que ele fosse dois: o texto + o meta-texto. Notas de rodapé robustas, como se o texto é que comentasse sobre elas, e não o contrário.

Iniciei este diário tendo em vista organizar minhas ideias para a escrita de um outro texto, que seria submetido para esta publicação. Logo notei como o próprio diário era o texto que eu queria escrever, no formato que mais me apetece escrever, e assim segui; sem, no entanto, abandonar a imagem mental de um outro texto do qual este seria um comentário, um meta-texto. Eis que surge o rodapé: o espaço perfeito no qual colocar o meu texto imaginário.

Escrevo, às avessas, notas de rodapé alçadas ao papel central, enquanto o protagonista vai para a coxia da página engrandecer-lhe. Comecei uns dois dias atrás a escrever pensando nas notas de rodapé que terei de fazer. Além de um diário de bordo, o que redijo é também uma carta de intenções\notaarquivo{notas/carta-de-intencoes.tex}, as quais venham a ser continuamente realizadas nos rodapés.

Imagino uma prática científica orientada pela e para a comunicação. Um método que parte do discorrer de uma carta de intenções bem localizada e bem endereçada. Depende dessa carta e do público ao qual se dirige a escolha de referenciais teóricos, métodos, dados e experimentos e demais mediadores do trabalho\notaarquivo{notas/localizacao-rodapes.tex}.

A escrita, e mais especificamente a escrita linear, possui limitações comunicativas sérias. A oralidade tem uma força ancestral de trazer a palavra encarnada no presente, enquanto a escrita deixa a palavra imóvel. Quero aqui propor uma forma de escrita malemolente e que se faça presente, para que a palavra escrita dance como a palavra cantada.

Não serve aqui um texto \enquote{objetivo}, com sentido único e bem definido. Isso é uma falácia do signo. Palavras se movem a cada uso, e, mediante sua ondulação semântica, vão emergindo dos textos novos sentidos.

Os rodapés estão aí para atualizações, correções e presentificações em geral. No método da carta de intenções, há esta vantagem de se poder ressituá-la em diferentes circunstâncias. Uma carta de intenções que aceite diferentes experimentos e teorias é mais articulada do que um artigo científico que só aceita replicações idênticas ao procedimento original.

A carta de intenções dá à ciência o poder espiralar da oralidade. No discurso oral, cada performance carrega o que se repete junto do que se diferencia localmente. Junto ao rodapé bem localizado, a carta de intenções faz da ciência uma arte declamatória; o rodapé  presentifica a performance.

O leitor fica sempre convidado a refazer ou alterar os rodapés conforme a ocasião; conforme sua própria especulação das intenções propostas. Em seu ensaio, Latour se indaga sobre os aromas: estariam nas moléculas químicas? Ou nas memórias que perfumam? Com práticas científicas orientadas para a comunicação, uma mesma carta de intenções sobre possíveis benefícios da aromaterapia, por exemplo, poderia ter seus rodapés escritos por cientistas das mais diversas áreas, ou mesmo pelos narizes profissionais da indústria perfumista.

Qualquer pessoa pode, em última instância, fazer ciência declarativa, um formato científico open source iterate-to-evaluate --- ITE\notaarquivo{notas/open-source-ite.tex}. Avalia-se uma carta de intenções iterando um novo texto de rodapé; avalia-se um texto de rodapé também com outro texto de rodapé. Com uma boa diagramação, uma carta pode até ser publicada com mais de um conjunto de textos de rodapé.

%%                  %%
%%                  %%
%%  22 de outubro   %%
%%                  %%
%%                  %%
\entrada{22 de outubro de 2026}{Cultura com aspas}

Todos os dias, antes de vir escrever aqui, tenho feito alguma leitura para me aquecer a mente e começar a movimentar as palavras -- os corpos -- que a compõem. Hoje, dei continuidade ao ensaio-crônica de Manuela Carneiro da Cunha já mencionado. Quero, agora, pensar com Manuela o \enquote{Antropoceno} com as aspas.

Em seu texto, a antropóloga distingue a cultura da \enquote{cultura} enquadrando-as em diferentes níveis lógicos. A segunda aparece como uma performance da primeira, e ainda, um metacomentário a seu respeito. A \enquote{cultura} como unidade elementar de um contexto interétnico composto por diferentes \enquote{culturas}; e a cultura como um agregado orgânico de práticas, costumes e saberes que organizam as diferenças (e são por elas organizadas) internamente a um coletivo de pessoas.

A \enquote{cultura} é uma performance cultural daquele coletivo frente a outros coletivos, constituindo uma rede de relações, não entre pessoas, mas entre grupos, entre clãs. E naquele contexto inter-\enquote{cultural}, indivíduos são as \enquote{culturas} elas próprias, que se distribuem homogeneamente entre as pessoas que as compõem. Assim, na Copa do Mundo, por exemplo, surge o \enquote{brasileiro}, e cada brasileiro real compartilha com os demais a experiência da exibição de nossa performance \enquote{cultural} perante outras \enquote{culturas}, e se aproxima subjetivamente desse \enquote{brasileiro} que vai se produzindo, se realizando, em suas performances; hoje, a três dias de uma votação que, independentemente de seu resultado, será um marco na história do Brasil, o país se divide entre dois clãs, numa eleição presidencial acirradíssima que pode constituir ordenamentos sociais profundamente distintos conforme o vencedor do pleito; na Copa, diante do Outro de fora, somos mais iguais; nas eleições, diante do Outro de dentro, somos mais diferentes; esta é a diferença entre estar no nível da \enquote{cultura} ou da cultura.

Quando proponho utilizar as aspas em \enquote{Antropoceno}, quero deslocar o discurso para o nível lógico de um contexto interétnico. Latour faz esse deslocamento em "Jamais fomos modernos" ao invocar os modernos e os não-modernos para a narrativa que tece; já em \enquote{\textit{How to talk about the body?}}, ele opera de outro modo, focalizando não o corpo em si, mas \guillemetleft falas sobre o corpo\guillemetright, ou seja, falando de um \enquote{corpo} com aspas. Quero mostrar a simetria entre o que fazem Bruno, Manuela, Ailton e Antônio, para então pensar o antropoceno a partir dessa simetria.

Sem aspas e com iniciação minúscula, o antropoceno é o tempo-espaço no qual vivem os antropocênicos; com aspas e em maiúscula, o \enquote{Antropoceno} se refere a \guillemetleft falas sobre o antropoceno\guillemetright. Minha proposta é de um \enquote{Antropoceno} no singular-processual que articule os diferentes discursos sobre o antropoceno. Proponho ainda que o \enquote{Antropoceno} se articula à patologia da comunicação que esconde a forma e apresenta os orientes como objetos.

Essa escolha terminológica denuncia, a despeito da significação ou datação que se articule a ela, o pressuposto que é a grande divisão entre sujeito e \enquote{objeto} que a Constituição moderna operacionaliza. Penso o \enquote{Antropoceno} como a epítome desta Constituição. Se vivemos um período de múltiplas crises planetárias, é porque vemos problemas crescerem e se avolumarem à nossa frente e não encontramos agência para atravessá-los; é porque enquanto antropocênicos estamos subtraindo uma infinidade de agenciamentos possíveis e capazes de nos projetar para um exo-antropoceno, e então para um pós-antropoceno.

Sair do antropoceno demanda saltar do regime da cultura para o da \enquote{cultura}. Com esse deslocamento lógico-semiótico, enxerga-se o \enquote{Antropoceno} e se abre um imenso leque de oportunidades para ficar com o problema do antropoceno. Embarca-se, assim, no exo-antropoceno, um navio interétnico no qual não-antropocênicos fazem a travessia para o pós-antropoceno, quando os antropocênicos estarão extintos.

O exo-antropoceno é o tempo-espaço quando não-antropocênicos constituem laços, saberes e práticas afetivas -- uma cultura exo-antropocênica -- para se imunizarem contra o \enquote{Antropoceno}. Ao longo da viagem, por vezes encontram antropocênicos à deriva no mar, sem uma forma na qual se agarrarem, e os resgatam; mas sabem que muitos se afogarão antes de conseguir embarcar; sabem também que, seja pelo afogamento, seja por embarcar no exo-antropoceno, o destino dos antropocênicos é a extinção. Apenas então poder-se-á aterrar na terra firme do pós-antropoceno.

%%                  %%
%%                  %%
%%  23 de outubro   %%
%%                  %%
%%                  %%
\entrada{23 de outubro de 2026}{O piquenique cósmico}

Hoje li algo muito bonito do ensaio de Manuela. Ela falava sobre o conhecimento e seus diversos regimes. Disse que entre os povos amazônicos, a natureza é idiossincrática, e a cultura -- enquanto modo de conhecê-la, é universal; até mesmo os animais não-humanos compartilhariam com os humanos essa aculturação da natureza, apenas com seus referentes específicos\notaarquivo{notas/refs/cultura-cultura-universal.tex}.

O trecho já estava grifado, e certamente me ajudou a construir a imagem de um piquenique cósmico a orientar os encontros entre viventes. Imagino primeiramente um bosque -- essa natureza idiossincrática -- no qual os diferentes seres passeiam acompanhados de seus modos particulares -- seus orientes -- de reconhecer as coisas do bosque e selecionar o que coletar em suas bolsas; num segundo momento -- do piquenique em si -- um conjunto de seres se reúne numa clareira no centro do bosque: estende-se uma toalha no chão; todos, então, se sentam e compartilham suas coletas; reunidos, cada qual toma da toalha aquilo que lhes apetece, tudo que reconhecem como orientes para seus passeios pelo bosque. Penso em Antônio nos enquadrando como compartilhantes da terra; e num piquenique entre seres que compartilham seus diferentes modos de aculturar a terra, suas diferentes sensibilidades ao se relacionar com ela e com os demais compartilhantes.

Assim como não se pode conhecer a cultura do Outro imediatamente, apenas sua \enquote{cultura} -- performance da primeira num contexto intercultural --, um compartilhante não acompanha o modo de coletar dos demais, tendo acesso apenas ao resultado da colheita que estes compartilham consigo. O encontro com o Outro é sempre intercultural, daí a inacessibilidade à sua cultura, que fala através de performances específicas a cada encontro, constituindo máscaras -- tipo especial de oriente, espécie de persona que constela um conjunto de orientes (símbolos, práticas, técnicas, costumes) estabilizadas e continuamente re-estabilizadas nos encontros entre diferentes compartilhantes -- que realizam duplamente a comunicação entre eles. Por um lado, é mediante as máscaras que a comunicação acontece; por outro, elas são a própria comunalidade constituída em si entre diferentes compartilhantes que se comunicam.

No \enquote{Antropoceno}, a capacidade de constituir máscaras foi se esgotando com o monoteísmo cosmofóbico. Em vez de piqueniques, os antropocênicos fazem guerras tacando os frutos de suas colheitas uns nos outros, pois são incapazes de enxergar na colheita alheia senão um veneno ao seu próprio modo de coletar. E não é de todo equivocada a visão, pois os piqueniques cósmicos são potentes imunizantes contra a cosmofobia antropocênica.

Se começarem a sentar-se à toalha com o Outro, os seres antropocênicos logo estariam fazendo colheitas muito mais ricas, comunicando com os demais compartilhantes modos muito mais agradáveis e seguros de viver com a terra. Rapidamente se sintonizariam de volta ao ritmo desta provedora de tudo. Mas é preciso que vejam que o que tá parecendo veneno é vacina, que alcancem o exo-antropoceno para que enxerguem o \enquote{Antropoceno}.

%%                  %%
%%                  %%
%%  24 de outubro   %%
%%                  %%
%%                  %%
\entrada{24 de outubro de 2026}{O paradoxo da cosmofobia}

Como dizer que o que parece uma coisa é outra sem soar colonialista, sem cair em cosmofobia? O piquenique é um encontro para trocas cosmológicas entre diferentes, mas aqueles que só sabem jogar sua colheita nos outros e que se nega a todo e qualquer fruto alheio não têm lugar nesta mesa. Seria este um caso análogo ao do paradoxo da tolerância?

Este diz que, para o convívio entre diferentes, é preciso recusar a presença do diferente intolerante, aquele que não tolera diferenças. O piquenique cósmico é uma heterotopia especulativa fundada não  na tolerância, mas na amizade e no apreço pelo diferente. Não seria de todo equivocado chamar ainda de uma experiência de amor ao próximo.

No piquenique, quanto mais diferenças se agregam e compartilham, mais amplo se torna o mundo de todos os presentes; mais articulado e com mais agenciamentos possíveis. Penso que o que permite superar o paradoxo é o seguinte: senta-se à mesa com quem se senta à mesa, joga-se frutos a quem joga frutos; trata-se de reciprocidade -- com os cosmofóbicos também se aprende, seus frutos lançados também se aproveitam: eles ensinam, ao menos, uma imunização contra aquilo que quer destruir o seu modo de existir. Ao Outro que teme a toalha no chão não se nega um lugar no banquete; tampouco se deixa de tomar para si os frutos arremessados que se reconheçam interessantes, nem de arremessar-lhes de volta os frutos por si próprio colhidos: isto se faz mostrando que a transformação dos modos de viver é o segredo para manter estes mesmos modos, e para uma boa vida com a terra.

% === adicionar página de Bispo === %
Bispo conta que os povos indígenas ensinaram aos povos afroconfluentes que a ancestralidade que funcionava em sua terra nativa, da qual foram retirados à força, precisava mudar para funcionar do lado de cá do Atlântico. Houvessem então tradições afroconfluentes avessas à transformação, estas se perderam pois sua rigidez as enfraquece. No jogo cósmico da vida, a cosmofobia pode até render bons frutos localmente, no contexto da cultura, por algum período de tempo; no longo prazo, contudo, a terra seleciona os compartilhantes capazes de se transformarem junto a ela e os demais compartilhantes: aqueles que não temem o piquenique cósmico.

%%                  %%
%%                  %%
%%  25 de outubro   %%
%%                  %%
%%                  %%
\entrada{25 de outubro de 2026}{Os antropocênicos e suas filas}

No instante em que cheguei na escola onde voto, tocou o sinal dando o aviso de que os portões deviam ser abertos. Diante de um espaço amplo, as pessoas faziam uma fila que se desfez assim que a larga passagem se abriu. Uma mulher deparou-se com dois policiais e os saudou: \enquote{que bom que estão aqui! Hoje tá complicado!}; \enquote{quê isso!?}, respondeu cordial e sorridente um dos soldados, \enquote{tá tudo tranquilo!}.

Fui o primeiro de minha seção a votar; isso porque não fui contaminado pela patologia da ordenação linear que acometeu os enfileirados que encontrei ao chegar ali. Foi bastante curioso e um tanto engraçado ver tanta gente numa linha reta aguardando a abertura de um portal pelo qual passariam seis ou sete pessoas ao mesmo tempo tranquilamente. Como não souberam confluir com o território, ordenaram-se sinteticamente numa linha reta desnecessária, ignorantes de que \enquote{tá tranquilo} se espalhar organicamente pela passagem.

O saber sintético é assim: em prol de uma ordem pela ordem, condena à profanidade um amplo espaço de movimento e agência; porque não sabe confluir senão com uma noção abstraída de \enquote{ordem} que faz com que se façam filas desnecessárias por onde passariam multidões. O \enquote{Antropoceno} é, também, isto: esta patologia da mente linear. As múltiplas crises, como que uma só crise de confluência: os espaços de agência encontram-se interditados, profanados pela doença da linearidade -- a cosmofobia está matando os antropocênicos.

Sou otimista porque penso que a crise cosmofóbica vem se agudizando profundamente. Os entregadores de aplicativo conhecem bem a confluência dos cruzamentos frente a morosidade da linha reta. Se os carros se comportam como um grande rebanho sendo direcionado e redirecionado num complexo de currais, as motos são como cães confluindo com a terra e seus compartilhantes, correndo entre os bichos, se espalhando com o território.

Sensibilidades para a confluência vêm emergindo e abrindo espaços. Vejamos o que se deu na campanha eleitoral para o segundo turno destas eleições presidenciais de 2026: metade do Brasil esteve em grande parte mobilizada por uma urgência existencial para construir confluências de toda ordem. Fazer campanhas políticas demanda mentalidades politeístas a fim de agenciar-se com o Outro.

A vitória preliminar de Flávio Bolsonaro no primeiro turno das eleições foi um gatilho poderoso para a imunização contra a cosmofobia na esquerda brasileira. Buscou-se ocupar todos os espaços de agenciamento: desde estatísticas e dados de votação do TSE, até a articulação de aproximações cosmopolíticas com os mais diferentes segmentos em diversos eixos. Como falar com o agro, como falar com as mulheres, com os nêutros, a mobilização para dar caronas e assistências em geral para que as pessoas possam ir votar; num curto período de tempo, a busca urgente por \enquote{virar votos} fez com que a arbitragem do jogo fosse a capacidade para constituir comunalidades junto aos mais diferentes Outros, numa abertura cosmopolítica como pouco se vê entre os antropocênicos.

Em termos latourianos, os corpos brasileiros de esquerda tornaram-se bem mais articulados -- pois sensíveis a mais diferenças no mundo -- do que eram antes da corrida eleitoral. Muito se aprendeu a ouvir o Outro, a confluir com os territórios -- especialmente o digital --, a olhar para o local que medeia o global. E independentemente de qual seja o resultado final, toda essa sensibilização já serviu para oxigenar a imaginação política contrariamente à linha reta, no sentido das confluências.

Pensar o \enquote{Antropoceno} articuladamente a uma patologia da comunicação abre espaços de agenciamento mediante articulações e cruzamentos ontoepistêmicos, cosmopolíticos. Desde o exo-antropoceno é mais fácil evitar as filas desnecessárias dos antropocênicos e confluir com os territórios, a terra, e seus diferentes compartilhantes. Uma prática científica politeísta é capaz de constituir comunicações mais potentes com os territórios, com compartilhantes humanos e não-humanos; uma cosmopolítica ativa e atenta é capaz de constituir arranjamentos coletivos, e inter-coletivos, mais confluentes; uma aculturação politeísta do mundo é mais ágil na ginga necessária para sobreviver à queda do céu.

%%                  %%
%%                  %%
%%  26 de outubro   %%
%%                  %%
%%                  %%
\entrada{26 de outubro de 2026}{Um ponto de virada?}

Ontem foi um dia com o potencial de ser daqueles que encaminham os rumos históricos a renovadas direções. Pensar no sucesso da campanha pela eleição de Lula após a derrota para o candidato do fascismo no 1\textsuperscript{o} turno é, também, pensar no \enquote{Antropoceno}. Porque não foram lulistas, petistas, progressistas, comunistas e eleitores de esquerda de modo geral os únicos que colaboraram com seus votos e sua voz para a \enquote{vitória da democracia}: o esforço reuniu pessoas da direita, ex-eleitores e ex-apoiadores de Jair Bolsonaro, e até mesmo a jornalista e ex-deputada Joice Hasselmann --- que foi a líder de seu governo no Congresso Nacional em 2019\notaarquivo{notas/refs/hasselmann.tex}.

Diante da extensiva patologia da comunicação que medeia e significa o \enquote{Antropoceno}, as últimas três semanas serviram como um tratamento de choque para a aprendizagem, a construção e a reprodução de vocabulários e de modos de fala e de escuta mais potentes para a constituição de comunalidades --- ou seja, para a confluência comunicacional. Os diferentes grupos e pessoas de algum modo territorializados à esquerda tendem a ser, ironicamente, os mais adoecidos pela cosmofobia. É verdade que o fascismo se nutre dos nacionalismos, racismos e intolerâncias mais agúdos; mas ele sabe como ninguém, como a direita sabe, agregar diferenças em torno das comunalidades que deseja constituir.

Não se trata meramente de saber como falar, ou da vantagem política e comunicativa de já ter os referenciais de nossos imaginários e ideários marcados pelo status quo. O poder do capital e do dinheiro vem de sua profunda capacidade de se articularem e mediarem praticamente todas as experiências da vida de (quase) qualquer pessoa\notaarquivo{notas/pobreza-criada.tex}; o poder da comunicação de direita está na tomada dessa potência que atravessa todos os territórios e subjetividades atravessados por ela, partir de uma perspectiva que estabelece sem grandes esforços que ter mais dinheiro melhora a vida de qualquer um, e assim pode constituir uma legião de diferentes -- conservadores, reacionários, liberais, do povo mais pobre ao povo mais rico, nos campos e nas cidades, dos mais diferentes gêneros e orientações sexuais, de (quase) todos os tipos, enfim, de territórios e identidades -- que concordam que a liberdade para ganhar dinheiro é a única coisa que pode lhes pôr nas mãos a potencialidade de, por seu próprio esforço, melhorarem suas vidas e de suas famílias.

Nessa controvérsia, a esquerda tem a desvantagem de propor e lutar por contrapontos, amortizadores, quando não pela destituição completa dessa força fundamental para a vida digna em virtualmente todos os territórios ocidentalizados. Mesmo em seus próprios territórios, o dinheiro e o capital aparecem como forças fundamentais; mas a comunicação costuma ser falha pois não se escuta o território do Outro atentamente para compreender como as políticas, propostas e linguagens esquerdistas muitas vezes retiram as pessoas do único jogo que conhecem e que -- bem ou mal -- aprenderam a jogar para viver dignamente.

A centralidade da noção de \enquote{disputa}, de \enquote{competição}, também é outro grande entrave. É preciso aprender a negar a disputa como o único tabuleiro disponível para a dignidade. Do contrário, os contrapesos ao poder econômico e social como as políticas públicas de distribuição de renda, ou para as defesas das populações mais vulneráveis social e economicamente, serão traduzidas localmente nos territórios como privilégios e vantagens indevidas a classes específicas de jogadores enquadrados como concorrentes num mesmo jogo.

%%                  %%
%%                  %%
%%  06 de novembro  %%
%%                  %%
%%                  %%
\entrada{06 de novembro de 2026}{Atravessamentos}
Um diário é atravessado por muitas coisas -- inclusive interrupções imprevistas -- pois ele é bem localizado relativamente à vida de quem o escreve. Uma contingência externa agora chega ofertando confluência com os propósitos originais deste diário. A escrita para o dossiê \enquote{Pensar o Antropoceno} recebe em suas águas aquelas da apresentação no MultiTÃO Pesquisas\notaarquivo{notas/multitao-pesquisas.tex}, confluindo no vale de \enquote{Cecinepá: uma heterotopia especulativa para o mundo comum} -- minha contribuição para o evento.

Comecei a vislumbrar Cecinepá em minha primeira leitura do ensaio latouriano supracitado quando o autor imagina um \enquote{arquipélago de conexões heterárquicas} como organização da ciência\notaarquivo{notas/latour-arquipelago.tex}. As ilhas seriam cientistas, laboratórios, artigos ou resultados, articulados entre si e validados caso a caso. Isto ocorreria no sentido oposto ao das classificações objetivistas, totalizantes e reducionistas da prática científica ocidental.

Comecei a imaginar uma unidade de conhecimento como uma ilha no MAR, que \guillemetleft é\guillemetright em relação a outras ilhas, enquanto elemento e corpo --- simultaneamente --- de um arquipélago. Assim como o \enquote{corpo} latouriano é descrito pelas diferenças às quais é sensível, uma ilha (de saber) o é pelas ilhas às quais se articula. Cheguei então ao SAL como matéria de tais ilhas: Saber-Afeto-Laço, três componentes indissociáveis.

Retomando a ilustração do corpo que se transforma através do treinamento com o kit-de-odores, destaco como o corpo que aprende algo é o corpo que se sensibiliza a novas afetações. Seja o perfumista aprendiz, ou o químico, seus corpos aprendem ao aprenderem a articular novos afetos (a forma-sujeito) a outros conhecidos (as formas-orientes). Não há, portanto, saber que não seja afetivo, nem afeto que não constitua saber\notaarquivo{notas/saber-afeto.tex}.

Por outro lado, se o nariz incorpora o kit-de-odores como parâmetro para as fragrâncias que encontra no mundo -- e o químico, o cromatógrafo para as moléculas --, um laço é constituído na transformação do corpo-aprendiz em corpo-sensibilizado. A aprendizagem costura o corpo a novas diferenças no mundo, e ao artifício alfaiate que a realiza. Assim, também se nota que não há aprendizado sem constituição de novos laços, e vice-versa\notaarquivo{notas/saber-laco.tex}.

O SAL é o elemento fundamental das ilhas em Cecinepá. O movimento das marés junto ao calor do SOL vai gradativamente formando corpos articulados de SAL em meio ao MAR de contingências. Estes três componentes --- SAL, SOL e MAR --- é que demarcam a geografia e a dinâmica de Cecinepá\notaarquivo{notas/sal-sol-e-mar.tex}.

Digo que se trata de uma heterotopia especulativa para o mundo comum, e penso que essa descrição mobiliza muitos corpos estranhos. Estes são como os odores indistinguíveis ao perfumista que inicia seu treinamento. Buscarei então elaborar um kit-de-odores para que meus interlocutores possam coabitar Cecinepá também.

O que Latour chama de \enquote{mundo comum} eu chamei de pós-antropoceno. O francês denota que epistemologias são sempre políticas e estabelecem um mundo comum. Mas aquela que orienta a ciência ocidental restringe aos \enquote{cientistas} o poder de dizer quais \enquote{essências}, qualidades primárias, compõem este mundo\notaarquivo{notas/latour-epistemologia-politica.tex}.

Prefiro, neste ponto, voltar a Nego Bispo e dizer que Latour critica a cosmofobia do ocidente, e busca uma prática científica politeísta, para compor o multiverso que enxerga ser um mundo verdadeiramente comum. Cecinepá, neste sentido, funciona como um artifício de imunização contra a cosmofobia em prol da composição contínua de um mundo politeísta. Nos meus termos, Cecinepá é a jornada pelo exo-antropoceno.

%%                  %%
%%                  %%
%%  07 de novembro  %%
%%                  %%
%%                  %%
\entrada{07 de novembro de 2026}{Uma insistência}

Mas por que uma heterotopia? Tenho me questionado por esta escolha. Cecinepá poderia ser apresentado como um mundo, simplesmente, e isso facilitaria muito minha vida em falar a respeito; mas desde sua idealização, sempre insisti em tratar este espaço como uma heterotopia.

%%                  %%
%%                  %%
%%  08 de novembro  %%
%%                  %%
%%                  %%
\entrada{08 de novembro de 2026}{Uma heterotopia}

Embora a noção de heterotopia me fascine, sempre é um transtorno pra mim falar a respeito. De algum modo, no entanto, eu faço questão de continuar trazendo-a comigo. Penso que heterotopias são melhores orientes do que utopias na movimentação em direção a outros mundos.

A heterotopia é um espaço diferente de todos os outros espaços. Isso significa que há uma articulação específica de afetos, símbolos e performances numa heterotopia. O shopping center, o hospital, a escola, são exemplos de heterotopias: cada um carrega uma etiqueta, uma liturgia própria; e -- o que mais me interessa -- cada heterotopia promove uma atualização específica de certos corpos.

A heterotopia é, portanto, um artifício. Logo, ela deve ser composta por orientes e sujeitos. Num shopping, orientes são os funcionários, os vendedores, as lojas e suas marcas; e sujeitos são as mercadorias disponíveis; o corpo que se transforma é o do consumidor, que encontra um universo de coisas, formas-sujeito que lhe afetam, mas ele não pode articular exceto mediante a compra, que articula ao corpo a coisa comprada, que passa a ser um oriente daquele corpo.

Cecinepá é uma heterotopia, não uma utopia. Não é um mundo pra gente almejar construir, que fica no horizonte inalcançável, desconectado do aqui-agora, chamando em sua direção. É uma linha de fuga que a gente traça a partir do local, e que nos orienta a imaginar, delinear e mapear outros mundos -- mundos reais, não utópicos\notaarquivo{notas/linha-de-fuga.tex}.

Penso Cecinepá como uma ponte para um mundo comum politeísta desde o mundo universal cosmofóbico que os antropocênicos tentaram em vão impor. Em vez de imaginar como seria esse mundo comum, Cecinepá é composto tendo em vista imunizar corpos antropocênicos contra a cosmofobia, trazê-los para fora do "Antropoceno". Então, no exo-antropoceno, todos os não-antropocênicos juntos podem efetivamente imaginar e continuamente compor o mundo comum politeísta.

%%                  %%
%%                  %%
%%  09 de novembro  %%
%%                  %%
%%                  %%
\entrada{09 de novembro de 2026}{Kit-de-traduções}

Acho que posso dizer que penso Cecinepá como uma contra-colônia. Foucault menciona como as colônias foram reservas para o imaginário europeu\notaarquivo{notas/refs/heterotopias-colonias.tex}; vejo isso acontecendo de dois modos: primeiro, como espaço para imaginar e constituir novos arranjamentos sociais; segundo, como um encontro-descoberta-engendramento de um Outro nos povos nativos, e que, a contra-pêlo, engendrava o Eu e o espaço metropolitanos. Cecinepá almeja desfazer esse engendramento subjetivo-espacial colonial e colonialista de mão dupla; e o faz constituindo-se um espaço para mediações onto-cosmológicas, uma praça pública na qual realizar piqueniques cósmicos.

O colegiado de pajés de que trata Manuela Carneiro é um desses piqueniques. Penso que este é o caso para todo encontro interétnico consciente desta sua condição. Ali havia diferentes aldeias, associações civis e instituições que performavam diferentes unidades étnicas, diferentes \enquote{culturas}, e costurar as diferenças culturais num contrato reconhecível e reconhecido por todas essas unidades era o objetivo do encontro.

A principal reivindicação dos pajés no início da reunião era que suas práticas recebessem o mesmo suporte e financiamento do Estado brasileiro recebidos pela medicina ocidental, o que causou grande estranhamento entre os representantes do ministério da saúde ali presentes\notaarquivo{notas/manuela-reivindicacao-pajes.tex}. Em sua crônica, a antropóloga mostra como cultura e \enquote{cultura} se distinguem por pertencerem a níveis lógicos distintos; mas, num contexto interétnico, além de se atravessarem estes dois níveis lógicos circunscritos a cada unidade étnica, também se atravessam os espaços lógicos das diferentes \enquote{culturas} em contato. Estes atravessamentos geram tropeços de tradução, estranhamentos como aquele citado, especialmente -- como foi o caso -- quando alguém articula um elemento da \enquote{cultura} do outro em sua comunicação; é justamente para lidar com estes tropeços que antropólogos participam de tais negociações, como tradutores interétnicos que garantem a compreensão mútua dos termos dos contratos; Cecinepá funciona, neste sentido, como um kit-de-traduções, que transforma antropocênicos em tradutores interétnicos.

Cada ilha de SAL em Cecinepá se articula a outras ilhas como um arquipélago. Quando o chefe indígena diz, em língua portuguesa, que \enquote{\textit{honi}} não é \enquote{cultura}, isso causa estranhamento à antropóloga, porque em seus próprios arquipélagos, as ilhas de \textit{honi} e cultura se encontram. Cecinepá ajuda a imaginar -- e, assim, a realizar -- uma comunicação sensível às diferenças de níveis, ou espaços, ou redes, ou arquipélagos lógicos; voltando a Foucault, agora no prefácio de \enquote{As palavras e as coisas}\notaarquivo{notas/as-palavras-e-as-coisas.tex} (...)

%%                  %%
%%                  %%
%%  15 de novembro  %%
%%                  %%
%%                  %%
\entrada{15 de novembro de 2026}{Corpos de encontro}

No dia 13, anteontem, estava marcada minha supracitada apresentação. Entretanto, ela não aconteceu, pois não consegui elaborar uma fala de 15 minutos que me soasse suficiente para falar de Cecinepá. O tempo era curto, e a linearidade do discurso, um imenso entrave.

Esta publicação, caso venha a se confirmar, será a primeira aparição em público deste projeto. Aqui faço uso das vantagens da palavra escrita, do jogo espiralar entre texto e rodapé, e de uma extensão virtualmente livre do meu discurso. Mas o fracasso com a minha apresentação me dispõe a terminar o diário rematando, numa única e final entrada, a apresentação de Cecinepá como uma heterotopia especulativa para o mundo comum, e o modo como este projeto propõe pensarmos o antropoceno.

Retomo o final da primeira entrada, quando relatei a mudança de rumos da pesquisa ao me encontrar com minha primeira orientação e o pensamento de Bruno Latour. Até aquele momento, eu buscava falar sobre o autismo nas veredas da ciência moderna. A partir de então, passei a buscar um outro modo de fazer ciência.

A ideia de Cecinepá surgiu e se tornou projeto quando, logo em seguida, me deparei com o edital para o mestrado no qual a elaboro hoje, e com o trabalho da professora Susana Oliveira Dias, orientadora de minha pesquisa, que cruza as fronteiras entre a arte e a ciência. Para aquele processo seletivo, logo notei a confluência entre as propostas da professora e o meu objetivo, e elaborei um projeto de pesquisa que a alimentasse. Assim, em sua origem, Cecinepá olha para o antropoceno articulando-o ao humanismo\notaarquivo{notas/humanismo.tex}, e propõe uma prática científica fabulativa para desarticular este olhar humanista que: de um lado, enquadra seres não-humanos como recursos desprovidos de agência; e de outro, limita as possibilidades de agência dos próprios humanos.

Em 20 meses corridos de pesquisa, notei oportunamente que era nas ocasiões em que precisei fazer comunicações a seu respeito -- em congressos e trabalhos para disciplinas -- que eu mais encontrava dificuldades; mas também era nestes momentos que eu mais avançava em minhas proposições. Como comunicar-se cientificamente se tornou, assim, sem que eu o buscasse, uma questão central para o meu trabalho; e que se articulava, de fato, à minha proposta de buscar um novo modo de fazer ciência. Logo me veio uma nova proposição, anterior e complementar à de Cecinepá, de pensar práticas científicas orientadas pela e para a comunicação.

Para cada oportunidade de me comunicar sobre minha pesquisa ao longo deste período, realizei uma pequena pesquisa diferente, orientada especialmente pelas circunstâncias da ocasião. Sempre se tratava, para mim, de dar uma forma a pensamentos e ideias informes em mim. A formatação, portanto, foi a própria realização científica; foi a articulação do que havia a se comunicar com as contingências de cada espaço comunicativo, isto é, de cada localização; e neste ato é que se fez ciência, no evento da comunicação, não antes dele.

Imaginar a comunicação como método vai ao encontro da noção de SAL --- saber-afeto-laço --- que se encontra em Cecinepá. Pois se todo saber envolve um laço, uma relação, não há como fazer ciência sem fazer comunicação. Desde este ponto de vista, veria a ciência moderna como uma prática orientada mais à comunicação com seus objetos, conceitos e categorias, por meio de experimentos e análises que os colocam no mesmo espaço comunicativo, do que à comunicação com outros humanos.

A escrita do presente texto, contudo, me levou a olhar para o artigo científico como uma forma discursiva orientadora de toda a prática científica hegemônica. Afinal de contas, trata-se do meio por excelência pelo qual se demarca a realização da ciência; é nele que se completa o ciclo epistemológico que atesta a construção de um novo saber mediante a aprovação pelos pares, certificada pela publicação em si do artigo. Não é de se espantar a centralidade recebida por tais publicações -- em quantidade e repercussão -- na avaliação da produção científica de pesquisadores na atualidade.

Pois estou aqui escrevendo um não-artigo, uma carta de intenções, uma forma que só me ocorreu em meio à sua própria escrita; e fazendo, assim defenderei, ciência. Pois foi neste ato de escrita, na formatação do SAL que carrego em meu corpo para as circunstâncias da publicação à qual escrevo, que constituí a \guillemetleft carta de intenções\guillemetright como um novo membro deste meu corpo, e que ofereço em comunicação. Corpos sensibilizados a tal forma habitam um mundo sensível diferenciado daquele em que a ciência somente se orienta pela forma-artigo; seus imaginários permitem pensar práticas científicas muito mais bem localizadas, capazes ainda de articular diferentes localizações, como disciplinas, teorias, ou qualquer outro conjunto de circunstâncias que constitua um locus de comunicação; portanto, de prática científica.

Assim como a carta de intenções, Cecinepá também é uma formatação para práticas científicas orientadas pela e para comunicação. Com esta heterotopia composta de MAR, SAL e SOL, fabula-se uma cosmologia capaz de mapear e articular diferentes cosmologias, ontologias e corpos. Logo, capaz de comunicar e pôr em relação diferentes epistemologias, modos de ser, de conhecer, de narrar o mundo; capaz, enfim, de especular.

Em Cecinepá, encontro meu corpo como um arquipélago composto por todas as ilhas de SAL que guardam tudo que sei, tudo que me afeta, e tudo com que me relaciono. A \guillemetleft carta de intenções\guillemetright é o seu próprio corpo, mas é também um membro do meu. Enquanto corpo, enxergo-a em suas duas faces: uma Oriental, na qual se encontram todas as ilhas de SAL que compartilhamos, sua constituição enquanto membro de meu corpo; e a face Outra, o MAR contingente e misterioso do qual podem emergir outras ilhas conforme as comunicações com o arquipélago em Cecinepá.

Seria possível imaginar diferentes métodos para mapear arquipélagos, corpos e membros, mas a força primordial de Cecinepá está na sensibilização dos imaginários humanos para as diferenças entre um corpo em si, e um corpo enquanto membro de outro corpo. O primeiro tem uma extensão, um conjunto de ilhas, inalcançável em sua totalidade a qualquer outro corpo; já o segundo tem uma extensão específica a cada corpo, e alcançável somente pelo corpo de que faz parte. Em outras palavras: Cecinepá permite imaginar cartograficamente, por meio de ilhas e arquipélagos, a diferença entre tudo que a \guillemetleft carta de intenções\guillemetright pode significar, e o que ela significa a cada corpo, humano ou não, sensível a ela -- a cada corpo que encontre nela uma face Oriental\notaarquivo{notas/dif-faces.tex}.

Vamos imaginar o \guillemetleft odor\guillemetright como um corpo, um \guillemetleft perfumista\guillemetright e um \guillemetleft químico\guillemetright como dois corpos dos quais o primeiro é um membro. Enquanto corpo, \guillemetleft odor\guillemetright se estende pelas ilhas de todas as faces Orientais que os diferentes corpos encontram nele, em meio ao MAR de sua face Outra; enquanto membro, ele é um subconjunto específico de ilhas para o \guillemetleft perfumista\guillemetright, e outro para o \guillemetleft químico\guillemetright -- cada um diante de uma face Oriental específica. Assim, numa comunicação hipotética entre estes dois, é possível que ambos se refiram ao \guillemetleft odor\guillemetright sem que se perguntem sobre qual seria o odor \enquote{verdadeiro}; o \guillemetleft odor\guillemetright enquanto corpo se torna uma ponte entre as sensibilidades de um e de outro, permite que se defrontem com a diferença entre o \guillemetleft odor\guillemetright para si e para o outro, e permite que se faça ciência na comunicação dessa diferença -- seja no sentido de buscar diminuí-la, seja no de articular as faces Orientais do \guillemetleft odor\guillemetright de ambos o \guillemetleft perfumista\guillemetright e o \guillemetleft químico\guillemetright a fim de constituir ou encontrar outras ilhas do \guillemetleft odor\guillemetright enquanto corpo\notaarquivo{notas/prova-do-odor.tex}.

Mais do que uma ciência orientada pela comunicação, Cecinepá permite imaginar e realizar práticas científicas diferenciais e mais bem localizadas. Bem localizadas porque localiza espacialmente os diferentes significados daquilo de que se fala. Diferenciais porque orientadas não em direção a uma Verdade sobre o mundo, mas em função das diferenças pelas quais o mundo se apresenta e constitui cada corpo; e ciências diferenciais abrem espaço a novos saberes, relações e afetações; ou seja, a novos agenciamentos que possam confluir com a vida diante das múltiplas crises do antropoceno.

Em Cecinepá, o \enquote{Antropoceno} é apenas \guillemetleft antropoceno\guillemetright, um corpo de encontro entre os diferentes discursos, nomeações e significações do antropoceno. Quando penso no \enquote{Antropoceno}, penso nos diversos corpos antropocênicos que dão origem e sustentação a este arquipélago. Quando penso no antropoceno sem as aspas, penso no tempo-espaço em que corpos antropocênicos se fazem presentes.

Os \guillemetleft antropocênicos\guillemetright buscam a Verdade. Buscam igualar o \guillemetleft antropoceno\guillemetright corpo ao \guillemetleft antropoceno\guillemetright membro nos diferentes corpos, mesmo que isso lhes custe submergir uma infinidade de corpos, saberes, afetos e laços. Pensar o \guillemetleft antropoceno\guillemetright com Cecinepá, por outro lado, desvia o pensamento de uma navegação gulosa pela igualdade entre corpo e membro, compromissada em dissolver ilhas, para uma navegação que se orienta pelas diferenças a fim de fazer --- buscando a confluência entre diferentes corpos --- com que se multipliquem as ilhas, as relações, os afetos, os agenciamentos.




